\chapter*{Abstract}

Urban mobility is inherently non-stationary: shifts in weather, demand, and network conditions alter data distributions and degrade model accuracy over time.
This diploma thesis designs, implements, and evaluates a drift-aware platform for continuous machine learning in Cooperative, Connected and Automated Mobility (CCAM).
The system combines a modular microservice architecture for model serving, online monitoring, multi-detector drift consensus, and automated retraining with hot-swap deployment.
A reproducible SUMO-based dataset for central Athens is constructed under normal and adverse-weather scenarios.
Concept drift is induced via friction reduction to emulate rain, preserving network topology while yielding measurable shifts (average speed $-13.68\%$, average trip duration $+16.74\%$).

Evaluation centers on Estimated Time of Arrival (ETA) prediction with a LightGBM model enriched by domain-specific features, exercised in a time-lapse experiment with an abrupt drift at the midpoint, due to rain conditions.
Under drift, baseline errors increase markedly (MAE $30.36\,s \rightarrow 56.93$\,s and MAPE $13.20\% \rightarrow 19.02\%$).
After detecting drift, the model is retrained using one hour of post-drift data and hot-swapped, recovering performance, when compared to the baseline model (MAE $-11.00$\,s, $-25.42\%$ and MAPE $-2.34$\,pp, $-13.85\%$).
The platform's flexible architecture is further evidenced by the evaluation of Fuel Consumption and Number of Stops models, while its dashboard and user interface provide real-time interpretability.
Overall, the results demonstrate a practical, closed-loop approach to detecting, responding to, and mitigating concept drift in realistic urban traffic settings.

\textbf{\textit{Keywords:}}
Continuous Machine Learning, Concept Drift, ETA Prediction, SUMO, MLOps, Gradient Boosted Trees, Cooperative Connected and Automated Mobility, Urban Mobility

\cleardoublepage
